\documentclass[11pt]{article}
\usepackage[T1]{fontenc}
\usepackage[margin=1in]{geometry}
\title{A Walking Tour of the Old Observatory}
\author{FlashTeX Fixtures Team}
\date{}
\begin{document}
\maketitle
\section{The entrance hall}
\label{sec:hall}
You are on page~\thepage\ of this tour. The hall clock has kept sidereal time since 1893.\footnote{The clock stopped for eleven days in the winter of 1947, when the dome froze shut and nobody could reach the winding mechanism.} Beyond the hall lies the instrument gallery of Section~\ref{sec:gallery} on page~\pageref{sec:gallery}, and the tour ends at the dome of Section~\ref{sec:dome} on page~\pageref{sec:dome}. The brass plaque halfway up the stairs lists every director of the observatory, including the founder, whose telescope still points through the dome slit. School groups gather here at nine, and the guides divide them into parties of twelve before anyone is allowed upstairs. The floor is original tile, worn into shallow paths between the door, the stairs, and the clock, and the trustees refuse to replace it because the wear pattern is itself a record of a century of visitors. Photography is allowed in the hall but not upstairs, a rule dating to the flash-powder era that nobody has bothered to repeal. The guest book on the stand by the door holds signatures going back to 1903, including three astronauts and one king, or so the guides claim when the queue moves slowly. Parties too large for one group wait here among the display cases, which hold spare eyepieces, a cracked objective from the dropped refractor of 1931, and the notebook in which the founder recorded the first season of observations. The cases are locked but well lit, and the small cards beside each object were handwritten by the curator in 1974 and have never been updated, so the prices mentioned on them are a quiet joke among returning visitors.
\section{The instrument gallery}
\label{sec:gallery}
The gallery (reached from the hall of Section~\ref{sec:hall}) holds the meridian circle and the blink comparator.\footnote{The comparator arrived in 1921 in a crate labelled ``glass, extremely fragile'', and unpacking took four staff members the better part of a day; it worked on the first evening, which the logbook records with an uncharacteristic exclamation mark. The meridian circle is older: it was ordered in 1875, delayed by a dispute over the mounting, delivered in 1879, and calibrated against the transit of Mercury the following spring. Both instruments were used nightly until 1963, when the new hill station took over the survey work; since then they have been kept in working order for demonstrations, and the comparator still blinks pairs of archival plates for visitors every hour.} Every dial here is original, and the wall charts show the proper motions measured on this very floor. The gallery is long and narrow, with cabinets of eyepieces along one wall and the plate archive along the other, and it takes a full quarter hour to walk its length while reading the labels. Evening lectures are held here in winter, when the dome is too cold for groups. The lecture series began in 1952 with a talk on the canals of Mars, and the posters for every talk since are pasted in overlapping layers on the back of the door, a collage of changing fashions in astronomy. Attendance is free but the room seats only forty, so regulars arrive early and the latecomers listen from the stairs, which the fire code tolerates because the building predates it.\footnote{The logbooks for both instruments fill an entire cabinet, and the handwriting changes abruptly in 1933, when the night assistant retired and his successor took over mid-volume with a steel nib instead of a quill; researchers date undated pages by that nib mark alone. The lecture posters are archived the same way, one volume per decade, and the 1960s volume is twice as thick as the others because the curator saved every duplicate, convinced the misprints would one day be valuable.}
\section{The dome}
\label{sec:dome}
The dome caps the tour that began on page~\pageref{sec:hall} in Section~\ref{sec:hall} and passed the gallery of Section~\ref{sec:gallery} on page~\pageref{sec:gallery}. The shutter still opens by hand crank,\footnote{Turning the crank takes a full minute of steady effort; children ask whether the whole sky moves, and the guides answer that it does, slowly.} and on clear nights the founder's telescope shows Saturn's rings. The tour ends here, two floors above the hall where it started, and parties are led back down the same stairs past the plaque of directors. On the way down the guides take questions, the most common being whether the telescope can see the landing sites on the Moon, to which the answer is a practiced and gentle no.
\end{document}
